\documentclass[tikz,border=2mm]{standalone}
\usepackage{chemfig}
\setchemfig{atom sep=2em}
\begin{document}
\schemestart \subscheme{ \chemname{\chemfig{C(-[3]CH_3)(-[5]H)=C(-[1]CH_3)(-[7]H)}}{cis-but-2-eno} \hspace{2cm} \chemname{\chemfig{C(-[3]CH_3)(-[5]H)=C(-[7]CH_3)(-[1]H)}}{trans-but-2-eno} } \schemestop --> ![...](imagenes/tema08/cis.svg){style="display: block; margin: 0 auto; width: 50%;"} * **Isomería Óptica:** Ocurre ante la presencia de un carbono asimétrico o quiral (unido a 4 sustituyentes distintos). Da lugar a enantiómeros, imágenes especulares no superponibles que desvían la luz polarizada hacia la derecha (dextrógiro) o hacia la izquierda (levógiro). <!-- ##chemfig id=optica sep=2em \schemestart \subscheme{ \chemname{\chemfig{CH_3-CH_2-C(-[6]Cl)(-[2]OH)-CH_3}}{2-clorobutan-2-ol} \hspace{2cm} \chemname{\chemfig{[,1.7] C(-[2]H)(-[6]I)(<[:330]Cl)(<:[:210]Br)} }{\footnotesize bromocloroyodometano} } \schemestop --> ![...](imagenes/tema08/optica.svg){style="display: block; margin: 0 auto; width: 60%;"} ## **5. Reacciones orgánicas principales** **Ruptura de enlaces e intermedios de reacción** Cualquier reacción química consiste en la ruptura de unos enlaces y la formación de otros nuevos. Según como se rompan podemos tener: **1. Ruptura homolítica u homopolar**: cada átomo participante del enlace retiene un electrón del par que constituía la unión. ($\ce{H3C-CH3 \; \rightarrow \; H3C· + ·CH3}$) **2. Ruptura heterolítica o heteropolar**: los dos electrones que constituyen el enlace son asignados al mismo fragmento. $\ce{A : B \; \rightarrow \; A^+ + B^−}$ La primera da lugar a radicales libres, la segunda a carbocationes o carbaniones. **Principales reactivos orgánicos** - **Radicales libres** - **Nucleófilos** (“amantes de lo positivo”, eléctricamente hablando) - **Electrófilos** (“amantes de lo negativo”) **Tipos de reacciones** **1. Reacciones de Sustitución** Un átomo o grupo de átomos de la molécula sustrato es sustituido por otro. a) **Reacciones de sustitucion nucleófila** Se producen cuando un reactivo nucleófilo, Y, sustituye a un átomo o grupo de átomos electronegativo, X, que está unido a un carbono: $\ce{R-X + Y \; \rightarrow \; R-Y + X}$ Las reacciones típicas son: - **Halogenación de alcanos:** $\ce{CH4 + Cl2 ->[luz] CH3Cl + HCl}$ - En **derivados halogenados formación de alcoholes**: el reactivo es el ión $\ce{OH^-}$ por ejemplo del NaOH: se cambia el halógeno por el OH y se obtiene un alcohol: $\ce{CH3-CH2-Br + NaOH \; \rightarrow \; CH3-CH2-OH + NaBr}$ - A partir de derivados halogenados también se pueden obtener otras sustancias: aminas, nitrilos, éteres, ésteres, pero la más importante es la que da alcoholes. - **Formación de derivados halogenados a partir de alcoholes**. Es justamente la reacción contraria a la anterior. Están catalizadas por ácidos. Si para la anterior se usaba un hidróxido para ésta al revés, se usa un ácido, por ejemplo: $\ce{CH3-CH2-OH + HCl \; \rightarrow \; CH3-CH2-Cl + H2O}$ b) **Reacciones de sustitución electrófila** Son típicas del benceno y los compuestos aromáticos. En ellas un hidrógeno del benceno se sustituye por el reactivo electrófilo. Ejemplos: <!-- ##chemfig id=bromobenceno sep=2em \schemestart \subscheme{ \chemname{\chemfig{*6(=-=-=-)}}{benceno} \; + \; \chemfig{Br_2} \; \arrow{->}[0,0.7] \chemname{\chemfig{*6(=-=(-Br)-=-)}}{bromobenceno} \; + \; \chemfig{HBr} } \schemestop --> ![...](imagenes/tema08/bromobenceno.svg){style="display: block; margin: 0 auto; width: 50%;"} <!-- ##chemfig id=nitrobenceno sep=2em \schemestart \subscheme{ \chemname{\chemfig{*6(=-=-=-)}}{benceno} \; + \; \chemfig{HNO_3} \; \arrow{->}[0,0.7] \chemname{\chemfig{*6(=-=(-NO_2)-=-)}}{nitrobenceno} \; + \; \chemfig{H_2O} } \schemestop --> ![...](imagenes/tema08/nitrobenceno.svg){style="display: block; margin: 0 auto; width: 50%;"} Las reacciones suelen estar catalizadas por ácidos. **2. Reacciones de Adición** Ocurren en moléculas insaturadas (dobles o triples enlaces). Los átomos del reactivo se añaden rompiendo el enlace $\pi$. También pueden ser por radicales libres, nucleófilas o electrófilas. Las más importantes son las electrófilas. Los ejemplos típicos son la **adición de hidrógeno** (H-H) (**hidrogenación**), de **halógeno** (ej: Cl-Cl), de **hidrácido** (ej: H-Br) o de **agua** (H-OH). Cuando el **reactivo es simétrico** no hay dudas sobre el compuesto obtenido, por ejemplo, la reacción de cloro (**simétrico**) con propeno da lugar a 1,2- dicloropropeno: $\ce{CH3-CH=CH2 + Cl2 \; \rightarrow \; CH3-CHCl-CH2-Cl + HCl}$ Sin embargo, cuando el **reactivo es asimétrico** podemos obtener dos isómeros, por ejemplo, con el mismo propeno, si reacciona con ácido bromhídrico podemos obtener el 2-bromopropano o el 1-bromopropano. <!-- ##chemfig id=propeno sep=2em \schemestart[0, 0.8, 1.5] \chemname{\chemfig{CH_3-CH=CH_2}}{propeno} \arrow{0}[,0.7] + \arrow{0}[,0.7] \chemname{\chemfig{HBr}}{ácido bromhídrico} \arrow{0}[,0.25] \arrow(.mid--.mid){->}[0,0.8] \arrow{0}[,0.75] \chemname{\chemfig{CH_3-CH(-[6]Br)-CH_3}}{mayoritario} \arrow{0}[,0.2] + \arrow{0}[,0.2] \chemname{\chemfig{CH_3-CH_2-CH_2(-[6]Br)}}{minoritario} \schemestop --> ![...](imagenes/tema08/propeno.svg){style="display: block; margin: 0 auto; width: 100%"} * Siguen la **Regla de Markovnikov**: El hidrógeno del reactivo se adiciona mayoritariamente al carbono insaturado que tenga ya más átomos de hidrógeno. * **Ejemplo (Hidratación de alquenos):** $\ce{CH3-CH=CH2 + H2O ->[H+] CH3-CH(OH)-CH3}$ *(Producto mayoritario: 2-propanol)* **3. Reacciones de Eliminación** Dos átomos o grupos de átomos se escinden de posiciones adyacentes para dar lugar a una insaturación (doble o triple enlace). * Siguen la **Regla de Saytzeff**: Se elimina preferentemente el hidrógeno del carbono vecino que posea menos hidrógenos, obteniéndose el alqueno más sustituido y estable. * **Ejemplo (Deshidratación de alcoholes):** $\ce{CH3-CH2-CH(OH)-CH3 ->[H2SO4][\Delta] CH3-CH=CH-CH3 + H2O}$ *(Producto mayoritario: but-2-eno)* **4. Reacciones de Oxidación y Reducción** * **Oxidación:** Un alcohol primario se oxida a aldehído y este consecutivamente a ácido carboxílico. Un alcohol secundario se oxida a cetona. * **Reducción:** Es el proceso inverso al anterior. **Problemas complementarios (Estilo PAU)** Ejercicio 1 Para cada uno de los siguientes procesos, formula la reacción, indica el nombre de los productos y el tipo de reacción orgánica: * **a) Hidrogenación catalítica del 3-metilbut-1-eno.** $\ce{CH3-CH(CH3)-CH=CH2 + H2 ->[Pt/Pd] CH3-CH(CH3)-CH2-CH3}$ *Producto:* 2-metilbutano. *Tipo:* Adición (Reducción). * **b) Deshidratación de 1-butanol con ácido sulfúrico en caliente.** $\ce{CH3-CH2-CH2-CH2OH ->[H2SO4][\Delta] CH3-CH2-CH=CH2 + H2O}$ *Producto:* but-1-eno. *Tipo:* Eliminación. * **c) Deshidrohalogenación de 2-bromobutano con hidróxido de potasio en etanol.** $\ce{CH3-CHBr-CH2-CH3 + KOH ->[etanol] CH3-CH=CH-CH3 + KBr + H2O}$ *Producto:* but-2-eno (mayoritario por la regla de Saytzeff). *Tipo:* Eliminación. * **d) Obtención de metanol a partir de clorometano.** $\ce{CH3Cl + NaOH -> CH3OH + NaCl}$ *Producto:* Metanol. *Tipo:* Sustitución nucleófila.
\end{document}
